\chapter{Introduzione al laboratorio}
\label{cap:introduzione}

\section{Che cosa faremo}

Questo laboratorio insegna a usare la programmazione matematica come strumento di
\emph{decisione}, non solo di calcolo. Ogni capitolo parte da un problema gestionale
concreto --- quanto produrre, dove localizzare un servizio, quale prezzo fissare, quanto
rischio accettare --- e lo trasforma in un modello di ottimizzazione che viene poi
implementato in Python, risolto con un solver e, soprattutto, \emph{interrogato}: quanto
vale un'ora di capacità in più? La soluzione resiste se i dati cambiano del 5\%? Quale
compromesso consiglieremmo a un decisore?

\begin{insight}[La domanda giusta]
Alla fine di ogni esercitazione la domanda non è soltanto \emph{``qual è l'ottimo?''},
ma \emph{``quale decisione suggeriamo e quanto è robusta?''}. Il valore di un modello
sta nella qualità della raccomandazione che permette di formulare.
\end{insight}

\section{Le tre classi di modelli e le sigle usate}

In tutta la dispensa i modelli appartengono a tre classi, indicate con le sigle
standard internazionali:

\begin{itemize}
  \item \textbf{LP} (\emph{Linear Programming}, programmazione lineare): obiettivo e
        vincoli sono funzioni lineari delle variabili;
  \item \textbf{QP} (\emph{Quadratic Programming}, programmazione quadratica):
        obiettivo quadratico, vincoli lineari;
  \item \textbf{NLP} (\emph{Nonlinear Programming}, programmazione non lineare):
        obiettivo o vincoli non lineari generali.
\end{itemize}

Un problema è \emph{convesso} quando ogni minimo locale è anche globale (definizione
precisa nel capitolo~\ref{cap:richiami}): per LP è sempre vero, per QP e NLP dipende
dalle funzioni, e lo segnaleremo ogni volta. Le altre sigle ricorrenti sono definite
al primo uso nel rispettivo capitolo.

\section{Notazione}

In tutta la dispensa vale la stessa convenzione:

\begin{itemize}
  \item \textbf{scalari e indici}: lettere minuscole in corsivo
        ($x_{it}$, $b_t$, $\lambda$, $i$, $t$);
  \item \textbf{vettori}: lettere minuscole in \textbf{grassetto}
        ($\vet x$, $\vet w$, $\boldsymbol\mu$); tutti i vettori sono colonna e
        $\vet a\T \vet b$ indica il prodotto scalare;
  \item \textbf{matrici}: lettere maiuscole in \textbf{grassetto}
        ($\mat Q$, $\mat A$); $\mat Q \succeq 0$ significa semidefinita positiva;
  \item \textbf{insiemi}: lettere maiuscole in corsivo ($I$, $T$, $S$, $N$, $A$);
        $|S|$ è la cardinalità.
\end{itemize}

Tre eccezioni, dichiarate per aderenza alla letteratura: le \emph{variabili
aleatorie} sono maiuscole ($D$ la domanda, $L$ la perdita) e le loro realizzazioni
minuscole ($d_s$, $\ell_s$); $F$ indica la funzione di ripartizione; il parametro di
regolarizzazione della SVM si scrive $C$ come in tutta la letteratura di machine
learning.

\section{Una scelta di campo: solo variabili continue}

Tutti i modelli di questa dispensa usano \textbf{esclusivamente variabili continue}:
niente variabili binarie o intere. È una scelta deliberata, con tre motivazioni.

\begin{enumerate}
  \item \textbf{Didattica}: con variabili continue valgono la dualità, i prezzi ombra,
        le condizioni KKT --- gli strumenti che trasformano una soluzione numerica in
        una spiegazione economica. Con variabili binarie questi strumenti si perdono
        o si indeboliscono.
  \item \textbf{Copertura}: una quantità sorprendente di decisioni reali è naturalmente
        continua: quantità, flussi, quote di portafoglio, prezzi, potenze, coordinate.
  \item \textbf{Progressione}: la stessa applicazione può passare dalla programmazione
        lineare (LP) a quella quadratica (QP) e non lineare (NLP), mostrando come cresce
        l'espressività del modello senza cambiare strumento.
\end{enumerate}

\section{Mappa dei modelli}

\begin{center}\small
\begin{tabular}{clll}
\toprule
Cap. & Applicazione & Classe & Concetto centrale \\
\midrule
\ref{cap:produzione} & Produzione e scorte multiperiodali & LP / QP & dualità, capacità, smoothing \\
\ref{cap:supplychain} & Supply chain e sostenibilità & LP / NLP & flussi, congestione, CO$_2$ \\
\ref{cap:markowitz} & Portafoglio di Markowitz & QP & rischio-rendimento \\
\ref{cap:pricing} & Pricing e revenue management & NLP & elasticità e capacità \\
\ref{cap:budget} & Budget pubblicitario & NLP conv. & rendimenti decrescenti \\
\ref{cap:localizzazione} & Localizzazione continua & NLP conv. & efficienza-equità \\
\ref{cap:ricaricaev} & Ricarica di veicoli elettrici & LP / QP & energia, picchi \\
\ref{cap:code} & Capacità e tempi di attesa & NLP conv. & capacità-attesa \\
\ref{cap:newsvendor} & Newsvendor e varianti & LP stoc. & decidere prima di sapere \\
\ref{cap:varcvar} & VaR e CVaR & LP & rischio di coda \\
\ref{cap:svm} & Support Vector Machine & QP & ottimizzazione per il ML \\
\bottomrule
\end{tabular}
\end{center}

\section{Struttura di ogni capitolo}

Ogni capitolo applicativo segue la stessa scaletta in otto passi, pensata per essere
replicata nelle relazioni di laboratorio:

\begin{enumerate}
  \item \textbf{Motivazione gestionale} --- il problema raccontato come decisione aziendale;
  \item \textbf{Costruzione guidata del modello} --- insiemi, parametri, variabili, vincoli
        e obiettivo, spiegati riga per riga;
  \item \textbf{Esempio numerico svolto a mano} --- un'istanza minima risolta con tutti i
        passaggi, per capire il meccanismo prima di accendere il solver;
  \item \textbf{Caso di studio} --- un'istanza realistica con dati riproducibili (CSV);
  \item \textbf{Implementazione Python/Gurobi} --- il codice completo, commentato;
  \item \textbf{Risultati} --- l'output del solver tradotto in linguaggio manageriale;
  \item \textbf{Analisi di sensitività} --- prezzi ombra, scenari, trade-off;
  \item \textbf{Esercizi} --- dalla verifica alla piccola estensione di ricerca.
\end{enumerate}

\begin{attenzione}
I numeri riportati in ogni capitolo sono generati dagli script della cartella
\texttt{python/}: eseguendo \texttt{python3 esegui\_tutti.py} si rigenerano dati,
tabelle e figure. Se modificate i dati, i numeri della dispensa vanno rigenerati.
\end{attenzione}

\section{Organizzazione del corso e valutazione}

Il percorso essenziale prevede quattro laboratori (capitolo~\ref{cap:organizzazione}
per il dettaglio): (1) produzione e scorte con dualità e prezzi ombra; (2) Markowitz e
frontiera efficiente; (3) modellazione non lineare (pricing o budget); (4) progetto a
scelta con presentazione manageriale. La valutazione suggerita pesa: correttezza della
formulazione 30\%, implementazione e verifica numerica 25\%, analisi di sensitività
25\%, interpretazione manageriale e qualità delle visualizzazioni 20\%.
